\uuid{WOkH}
\titre{Logique}
\auteur{bodin}
\organisation{exo7}
\licence{CC BY-NC-SA 4.0}
\datecreate{2020-01-21}
\niveau{L1}
\module{Algèbre}
\chapitre{Logique, ensemble, raisonnement}
\sousChapitre{Autre}
\difficulte{}
\theme{Logique}
\exo7qcmid{arnaud.arnaud.006}

\contenu{
\begin{question}
\qclassification{100.99}
Soit la suite $(x_n)_{n\in\Nn^*}$ définie par $x_n =
\frac{(-1)^n}{n}$.
\begin{answers}
    \bad{$\exists N > 0 \quad \forall n \in \Nn^* \qquad (n \ge N \implies x_n \ge 0)$.}

    \bad{$\exists \epsilon >0 \quad \forall n \in \Nn^* \qquad x_n \le \epsilon$.}

    \good{$\forall N \in \Nn^ * \quad \exists n \ge N \qquad  x_n < 0$.}

    \bad{$\exists n \in \Nn^* \qquad x_n =0$.}

    \good{$\forall \epsilon > 0\ \ \exists N \in \Nn^* \ \ \forall n \in \Nn^* \ \ (n\ge N \implies |x_n|\le \epsilon)$.}
\end{answers}
\end{question}
}
